% DRAIN OUT - Final Product Definition and Slide Content
% Compile on Overleaf with pdfLaTeX. No external files needed.

\documentclass[11pt,a4paper]{article}

\usepackage[margin=2cm]{geometry}
\usepackage[T1]{fontenc}
\usepackage{lmodern}
\usepackage{microtype}
\usepackage{amsmath}
\usepackage{booktabs,tabularx,array,longtable}
\usepackage{enumitem}
\usepackage[table]{xcolor}
\usepackage[most]{tcolorbox}
\usepackage{titlesec}
\usepackage{fancyhdr}
\usepackage[colorlinks=true,linkcolor=accent,urlcolor=accent]{hyperref}

% ---------- Look ----------
\definecolor{accent}{HTML}{0B5CAD}
\definecolor{accent2}{HTML}{E8750A}
\definecolor{soft}{HTML}{EEF3FA}
\definecolor{ink}{HTML}{1B2330}

\titleformat{\section}{\Large\bfseries\color{accent}}{\thesection}{0.8em}{}
\titleformat{\subsection}{\large\bfseries\color{ink}}{\thesubsection}{0.7em}{}

\setlength{\headheight}{14pt}
\pagestyle{fancy}
\fancyhf{}
\lhead{\small\color{accent}\textbf{DRAIN OUT} \textperiodcentered{} Final Product}
\rhead{\small Smart India Hackathon 2026}
\cfoot{\small\thepage}
\renewcommand{\headrulewidth}{0.4pt}

\setlist{itemsep=2pt,topsep=3pt}
\renewcommand{\arraystretch}{1.3}
\newcolumntype{L}{>{\raggedright\arraybackslash}X}

% A slide: title bar and the text that goes on it.
\newtcolorbox{slide}[1]{enhanced,breakable,colback=white,colframe=accent,
  coltitle=white,fonttitle=\bfseries\large,title={#1},arc=2pt,boxrule=0.8pt,
  left=8pt,right=8pt,top=6pt,bottom=6pt,before skip=10pt,after skip=10pt}
% A highlighted statement.
\newtcolorbox{keyline}{colback=soft,colframe=accent,boxrule=0.5pt,arc=2pt,
  left=8pt,right=8pt,top=5pt,bottom=5pt}
% A labelled block inside a slide.
\newcommand{\block}[1]{\par\smallskip{\color{accent2}\textbf{#1}}\par\smallskip}

% =====================================================================
\begin{document}

\begin{titlepage}
  \centering
  \vspace*{3cm}
  {\color{accent}\rule{\linewidth}{1.2pt}}\par\vspace{0.6cm}
  {\Huge\bfseries DRAIN OUT\par}
  \vspace{0.3cm}
  {\LARGE Final Product Definition\par}
  \vspace{0.4cm}
  {\large Street-Level Urban Flood Nowcasting \& Safe Routing\\ for Indian Cities\par}
  \vspace{0.6cm}
  {\color{accent}\rule{\linewidth}{1.2pt}}\par
  \vspace{1.5cm}
  \begin{tabular}{@{}ll@{}}
    \textbf{Event}     & Smart India Hackathon 2026 \\
    \textbf{Team}      & Tech Sprints \\
    \textbf{Prototype} & \url{https://github.com/QuantaBoy/SIH26085} \\
  \end{tabular}
  \vfill
\end{titlepage}

\tableofcontents
\newpage

% =====================================================================
\section{The Final Product}

\begin{keyline}
\textbf{DRAIN OUT} is a national platform that forecasts urban flooding
\textbf{street by street, 0--3 hours ahead, every 5 minutes}, detects overflowing
drains and manholes before they overflow, and guides citizens and emergency services
around the water. Any Indian city is onboarded from its own drainage records, without
new code.
\end{keyline}

\subsection{Product modules}

{\small
\begin{longtable}{@{}p{0.6cm}p{3.6cm}p{11.8cm}@{}}
\toprule
\textbf{\#} & \textbf{Module} & \textbf{What it delivers}\\
\midrule\endhead
1 & \textbf{Multi-source data ingestion} &
  IMD Doppler weather radar and rain gauges (rainfall); CartoDEM / LiDAR (terrain);
  city drainage records in GIS, CAD or scanned-map form; NRSC land use and soil maps;
  INCOIS tide levels for coastal cities; IoT water-level sensors in drains and
  underpasses; historical floods from NDMA, state disaster authorities and Sentinel-1
  satellite flood maps.\\
2 & \textbf{City onboarding engine} &
  Turns any city's drainage records into a working drain network: digitises paper maps,
  links drains into a network, infers flow direction from levels, auto-repairs missing
  links and unsized drains, and corrects terrain to surveyed levels.\\
3 & \textbf{Runoff model} &
  SCS-CN runoff for unpaved land from soil and land-use maps; imperviousness-based runoff
  for built-up areas.\\
4 & \textbf{Coupled hydraulic engine} &
  Drains and street surface simulated together: drain capacity (Manning), inlet capture,
  overflow and backflow at manholes, water spreading over streets towards outlets, tide
  level at coastal outfalls.\\
5 & \textbf{AI layer} &
  \textbf{ConvLSTM} trained on thousands of simulated storms and historical floods for
  city-wide forecasts in milliseconds; \textbf{Historical Pattern Model} for recurrence
  risk and new hotspots; \textbf{live correction} of the forecast from IoT sensor
  readings.\\
6 & \textbf{Risk and confidence} &
  Flood probability and severity class for every street, from repeated runs with the
  uncertain inputs varied.\\
7 & \textbf{Decision support} &
  Location-based early warnings; overflowing-manhole alerts; drain desilting priority
  lists; traffic diversion and evacuation suggestions.\\
8 & \textbf{Flood-safe routing} &
  A* and Dijkstra routes around streets forecast to flood during the trip, with modes
  for citizens, ambulances, fire services and traffic police.\\
9 & \textbf{Delivery} &
  GIS dashboard for city control rooms (Smart City ICCC); multilingual mobile app with
  push and SMS alerts; open REST APIs; alerts in the Common Alerting Protocol (CAP) for
  NDMA's SACHET system.\\
\bottomrule
\end{longtable}}

\subsection{Performance targets}

\begin{center}
\begin{tabular}{@{}ll@{}}
\toprule
\textbf{Metric} & \textbf{Target}\\
\midrule
Spatial resolution & Every street, flood depth in cm\\
Lead time & 0--3 hours, updated every 5 minutes\\
City-wide forecast refresh & Under 1 minute\\
Flood-safe route & Under 1 second\\
Accuracy & Each city validated against its past floods before go-live\\
Scale & One platform, many cities; each city is a data onboarding job\\
\bottomrule
\end{tabular}
\end{center}

\subsection{Users}
City corporations (urban local bodies) \textperiodcentered{} NDMA and state disaster
management authorities \textperiodcentered{} traffic police \textperiodcentered{} 108
ambulance and fire services \textperiodcentered{} citizens.

% =====================================================================
\section{How DRAIN OUT Addresses the Problem}

\begin{keyline}
\textbf{The problem.} Indian cities flood every monsoon, yet warnings name a city or a
ward, arrive with little lead time, and say nothing about which street, how deep or
when. Drains fail unseen until water is already on the road, desilting follows a fixed
calendar, and ambulances, buses and commuters drive into flooded roads and underpasses
because no one knows where the water will be.
\end{keyline}

\subsection{Problem to solution}

{\small
\begin{longtable}{@{}p{3.9cm}p{7.6cm}p{4.5cm}@{}}
\toprule
\textbf{Today's problem} & \textbf{How DRAIN OUT addresses it} & \textbf{Result}\\
\midrule\endhead
Warnings stop at city or ward level &
  Coupled drainage + terrain simulation computes flood depth on \textbf{every street},
  sampled every 20 m &
  ``Anna Salai, 35 cm, 4:40 PM'' instead of ``heavy rain in Chennai''\\
Little or no lead time; response is reactive &
  Rainfall nowcast drives the model \textbf{0--3 hours ahead, updated every 5 minutes} &
  Hours to act before water reaches the street\\
Drain failures are invisible until overflow &
  Every drain carries only its hydraulic capacity; the excess is traced to the
  \textbf{exact manhole} it overflows from &
  Overloaded drains and overflowing manholes known in advance\\
Desilting follows a fixed calendar &
  Drains ranked by how often and how badly they overflow across forecast storms &
  \textbf{Risk-based desilting}: crews go where drains will fail\\
Drainage records are incomplete, on paper or inconsistent &
  City onboarding engine digitises maps, links drains, infers flow direction and
  \textbf{auto-repairs missing links and unsized drains} &
  Any city can go live from the records it already has\\
Vehicles and ambulances drive into flooded roads &
  \textbf{A* / Dijkstra routing} costs every road by its forecast depth during the trip;
  roads over 30 cm are closed &
  Safe routes for citizens, 108, fire and police, in under a second\\
Recurring hotspots are not learned; climate shifts the risk &
  \textbf{Historical Pattern Model} learns each city's recurring flood zones and is
  refitted after every monsoon; streets flooding for the first time are flagged as
  \textbf{new hotspots} &
  Long-term planning targets the right places\\
Physics models are too slow for city-wide real time &
  \textbf{ConvLSTM trained on thousands of simulated storms} reproduces the physics in
  milliseconds; the physics engine stays the reference &
  City-wide refresh in under a minute, many cities at once\\
Forecasts carry no sense of certainty &
  Repeated runs with uncertain inputs varied give a \textbf{flood probability} per street &
  Officials know which warnings are firm\\
Departments work in silos &
  One live picture on a \textbf{GIS dashboard}, \textbf{open APIs} and \textbf{CAP alerts}
  to NDMA's SACHET &
  City corporation, police, hospitals and disaster authorities act on the same forecast\\
Citizens get generic alerts in one language &
  \textbf{Location-based, multilingual} push and SMS alerts for the user's street and route &
  Every citizen knows what it means for them\\
\bottomrule
\end{longtable}}

\subsection{From rainfall to action: one storm, hour by hour}

\begin{center}
\begin{tabular}{@{}p{2.2cm}p{13.4cm}@{}}
\toprule
\textbf{When} & \textbf{What DRAIN OUT does}\\
\midrule
\textbf{T\,--\,3 h} & Radar shows heavy rain approaching. DRAIN OUT forecasts which
  streets will flood, how deep and when, and which manholes will overflow.\\
\textbf{T\,--\,2 h} & Control room sees the drains at risk; maintenance crews are sent
  to clear the inlets and pumps that matter most.\\
\textbf{T\,--\,1 h} & Citizens on affected streets receive alerts in their language;
  traffic police set diversions; underpasses at risk are closed.\\
\textbf{T\,0} & Ambulances, fire services and commuters follow flood-safe routes; the
  forecast refreshes every 5 minutes as rain and sensor readings arrive.\\
\textbf{After} & Flooded streets are logged; the historical model and calibration are
  updated; drains that failed move up the desilting list for next season.\\
\bottomrule
\end{tabular}
\end{center}

\subsection{Why this works}
\begin{itemize}
  \item \textbf{It models the cause, not the symptom.} Street flooding happens when
        drains run out of capacity; DRAIN OUT simulates the drains themselves, so it
        knows where and when that happens.
  \item \textbf{Physics and AI check each other.} Physics gives trusted answers; AI gives
        speed and memory of past floods.
  \item \textbf{It ends in an action, not a map.} Every output is something a person can
        do: clean this drain, close this road, take this route.
\end{itemize}

% =====================================================================
\section{Slide Content}

% ---------------------------------------------------------------------
\begin{slide}{Slide 1 \textendash{} Title Page}
\begin{itemize}
  \item \textbf{Problem Statement ID:} \emph{(fill in)}
  \item \textbf{Problem Statement Title:} \emph{(fill in)}
  \item \textbf{Theme:} \emph{(fill in)}
  \item \textbf{PS Category:} Software
  \item \textbf{Team ID:} \emph{(fill in)}
  \item \textbf{Team Name:} Tech Sprints
\end{itemize}
\block{Subtitle}
DRAIN OUT: Street-Level Urban Flood Nowcasting \& Safe Routing for Indian Cities
\end{slide}

% ---------------------------------------------------------------------
\begin{slide}{Slide 2 \textendash{} Proposed Solution \& Unique Value Proposition}
\block{Proposed Solution}
\textbf{DRAIN OUT is a national platform that forecasts urban flooding street by street,
0--3 hours ahead, and guides people and emergency services around it.}
\begin{itemize}
  \item \textbf{Coupled drainage + terrain simulation} predicts flood depth in
        \textbf{cm on every street, every 5 minutes}.
  \item \textbf{Onboards any Indian city}: drainage maps (GIS, CAD, scanned) become a
        working drain network, with missing links auto-repaired.
  \item \textbf{Physics-trained AI (ConvLSTM)} with a \textbf{historical pattern model}:
        millisecond forecasts that learn each city's recurring flood zones.
  \item Detects \textbf{overflowing manholes and overloaded drains}, turning fixed
        desilting schedules into targeted ones.
  \item \textbf{Flood-safe routing (A* / Dijkstra)} for citizens, ambulances and traffic
        police.
  \item Delivered through a \textbf{control-room GIS dashboard, a multilingual mobile app
        and open APIs}, integrated with NDMA alerts.
\end{itemize}

\block{Unique Value Proposition}
\begin{tabularx}{\linewidth}{@{}LL@{}}
1.\ Street-level, not ward-level: depth in cm & 4.\ Manhole-level overflow and backflow detection\\
2.\ Missing-link-aware drainage network & 5.\ One platform, any Indian city\\
3.\ Physics + AI + historical memory & 6.\ Flood-safe routing in under a second\\
\end{tabularx}

\block{DRAIN OUT Addresses the Problem By}
\begin{tabularx}{\linewidth}{@{}p{5.2cm}L@{}}
\textbf{Street-level prediction} & Flood depth in cm on every street, 0--3 hours ahead, every 5 minutes\\
\textbf{Drain failure foresight} & Overloaded drains and overflowing manholes identified before they overflow\\
\textbf{Flood-aware rerouting} & Traffic, ambulances and fire services routed around flooded roads, reducing loss of life and congestion\\
\textbf{Targeted action} & Location-based multilingual alerts, and desilting sent where drains will fail\\
\end{tabularx}
\end{slide}

% ---------------------------------------------------------------------
\begin{slide}{Slide 3 \textendash{} Technical Approach}
\begin{enumerate}[label=\textbf{\arabic*.}]
  \item \textbf{Data collection:} IMD Doppler radar and rain gauges \textperiodcentered{}
        CartoDEM / LiDAR \textperiodcentered{} city drainage records \textperiodcentered{}
        NRSC land use and soil (SCS-CN) \textperiodcentered{} INCOIS tide levels
        \textperiodcentered{} historical floods (NDMA, Sentinel-1) \textperiodcentered{}
        IoT water-level sensors
  \item \textbf{City onboarding (preprocessing):} digitise $\rightarrow$ build drain
        network $\rightarrow$ auto-repair gaps $\rightarrow$ infer flow direction
        $\rightarrow$ correct terrain to surveyed levels
  \item \textbf{Physics layer:} SCS-CN runoff $\rightarrow$ drain capacity
        $Q_{cap} = \frac{1}{n} A R^{2/3} S^{1/2}$ $\rightarrow$ drains and street surface
        simulated together $\rightarrow$ overflow and backflow at manholes
  \item \textbf{AI layer:} ConvLSTM trained on simulated storms and historical floods
        \textperiodcentered{} Historical Pattern Model \textperiodcentered{} live correction
        from sensors \textperiodcentered{} flood probability per street
  \item \textbf{Decision support:} early warnings \textperiodcentered{} desilting priority
        \textperiodcentered{} diversions \textperiodcentered{} flood-safe routing
        (A* / Dijkstra)
  \item \textbf{Delivery:} GIS dashboard (ICCC) \textperiodcentered{} mobile app
        \textperiodcentered{} REST APIs \textperiodcentered{} NDMA SACHET (CAP) alerts
\end{enumerate}

\block{Outputs}
Street-level depth (cm) \textperiodcentered{} nowcast 0--3 h at 5-minute steps
\textperiodcentered{} flood risk and severity \textperiodcentered{} overflowing manholes
\textperiodcentered{} flood-safe routes

\block{Tech Stack}
Python \textperiodcentered{} FastAPI \textperiodcentered{} NumPy / SciPy
\textperiodcentered{} PyTorch (ConvLSTM) \textperiodcentered{} PostGIS
\textperiodcentered{} Leaflet / GIS \textperiodcentered{} Docker (cloud deployment)

\block{Prototype}
Working city-wide pilot on Chennai: \url{https://github.com/QuantaBoy/SIH26085}
\end{slide}

% ---------------------------------------------------------------------
\begin{slide}{Slide 4 \textendash{} Feasibility and Viability}
\block{Feasibility Analysis}
\begin{tabularx}{\linewidth}{@{}p{3.2cm}L@{}}
\toprule
\textbf{Data} & Every input has a national source: IMD, NRSC, Copernicus, INCOIS, NDMA, city drainage records.\\
\textbf{Modelling} & Built on established engineering methods: Manning, SCS-CN, CPHEEO, HEC-22.\\
\textbf{AI} & Trained on simulated storms, so it does not depend on scarce flood records.\\
\textbf{Scalability} & The same pipeline for every city; a new city is data onboarding, not new code.\\
\textbf{Integration} & Open APIs connect to Smart City command centres, 108, traffic police and NDMA alerts.\\
\bottomrule
\end{tabularx}

\block{Viability: Challenge $\rightarrow$ Mitigation}
\begin{tabularx}{\linewidth}{@{}p{4.4cm}L@{}}
\toprule
\textbf{Incomplete drainage records} & Auto-repair of missing links plus digitisation of paper maps\\
\textbf{Validation} & Each city tested against its past floods before go-live\\
\textbf{Model drift over years} & Recalibrated after every monsoon; historical model refitted on recent years\\
\textbf{Multi-agency integration} & Standard REST and CAP alert formats\\
\textbf{Nationwide scale} & Cloud deployment; one platform serving many cities\\
\bottomrule
\end{tabularx}

\block{Viability Statement}
DRAIN OUT is viable because it combines nationally available data with established
hydraulic engineering and AI in one modular pipeline that is validated on each city's
past floods and scaled city by city.
\end{slide}

% ---------------------------------------------------------------------
\begin{slide}{Slide 5 \textendash{} Impact and Benefits}
\begin{tabularx}{\linewidth}{@{}p{3.2cm}L@{}}
\toprule
\textbf{Safety} & Street-level warnings up to 3 hours ahead; flood-safe routes for citizens and emergency vehicles.\\
\textbf{Economic} & Less property and vehicle damage; desilting spent where drains actually overflow.\\
\textbf{Urban planning} & Recurring flood hotspots and undersized drains identified for long-term investment.\\
\textbf{Operational} & One shared picture for city corporations, disaster authorities, traffic police and hospitals.\\
\bottomrule
\end{tabularx}

\block{Core Impact}
Street-level early warning and safe mobility for India's flood-prone cities.

\block{Users}
City corporations \textperiodcentered{} NDMA / state disaster authorities
\textperiodcentered{} traffic police \textperiodcentered{} 108 ambulance and fire
services \textperiodcentered{} citizens

\block{Sustainable Development Goals}
SDG 3 Good Health \textperiodcentered{} SDG 9 Industry \& Infrastructure
\textperiodcentered{} SDG 11 Sustainable Cities \textperiodcentered{} SDG 13 Climate
Action
\end{slide}

% ---------------------------------------------------------------------
\begin{slide}{Slide 6 \textendash{} Research and References}
{\small
\begin{tabularx}{\linewidth}{@{}p{3.3cm}p{4.2cm}L@{}}
\toprule
\textbf{Source} & \textbf{Work} & \textbf{Gap DRAIN OUT fills}\\
\midrule
Ministry of Earth Sciences & I-FLOWS Mumbai & Ward-level alerts $\rightarrow$ street-level depth in cm\\
NCCR & C-FLOWS Chennai & Precomputed scenarios $\rightarrow$ live simulation every 5 min\\
NDMA & Urban Flooding Guidelines (2010) & Guidelines $\rightarrow$ an operational forecasting tool\\
MoHUA / CPHEEO & Storm Water Drainage Manual (2019) & Design standard $\rightarrow$ real-time drain load\\
FHWA & HEC-22 Urban Drainage Design & Inlet capture applied to live forecasting\\
IMD & Doppler Weather Radar network & Radar rainfall $\rightarrow$ street-level flood depth\\
\bottomrule
\end{tabularx}}

\smallskip
\textbf{Note:} put a source under every chart, and replace every ``LINK'' with the real URL.
\end{slide}

\end{document}
